% !TEX program = xelatex
% ============================================================
%  Arabic Kanban Board Template — Nibras Center
%  مشروع بحثي: لوحة كانبان
%  Columns: To Do | In Progress | Review | Done
%  Compile with: xelatex main.tex (twice)
% ============================================================

\documentclass[11pt,a4paper,landscape]{article}

\usepackage[a4paper,landscape,margin=1.5cm,top=2cm,bottom=2cm,headheight=12pt]{geometry}
\usepackage{array}
\usepackage{booktabs}
\usepackage{tabularx}
\usepackage{enumitem}
\usepackage{float}
\usepackage{titlesec}
\usepackage{fancyhdr}
\usepackage[table]{xcolor}
\usepackage{tikz}
\usetikzlibrary{positioning,calc,shapes.geometric}

\usepackage{bidi}
\usepackage[no-math]{fontspec}
\usepackage[quiet]{polyglossia}

\setmainfont[Scale=1]{NotoKufiArabic-Regular.ttf}
\newfontfamily\arabicfont[Script=Arabic,Scale=0.95]{NotoKufiArabic-Regular.ttf}
\newfontfamily\englishfont[Scale=0.95]{TeX Gyre Termes}

\setdefaultlanguage[calendar=gregorian,hijricorrection=1]{arabic}
\setotherlanguage[variant=british]{english}

\usepackage[breaklinks,hidelinks]{hyperref}
\hypersetup{pdftitle={لوحة كانبان},pdfauthor={الباحث الرئيسي}}

\pagestyle{fancy}
\fancyhf{}
\fancyhead[R]{\small\itshape لوحة كانبان}
\fancyhead[L]{\small اسم المشروع}
\fancyfoot[C]{\small\thepage}
\renewcommand{\headrulewidth}{0.3pt}

\definecolor{todoColor}{HTML}{E8E8E8}
\definecolor{progressColor}{HTML}{D4A017}
\definecolor{reviewColor}{HTML}{2E6DA4}
\definecolor{doneColor}{HTML}{27AE60}
\definecolor{headerColor}{HTML}{1A3C6B}

\newcolumntype{L}[1]{>{\raggedright\arraybackslash}p{#1}}
\newcolumntype{C}[1]{>{\centering\arraybackslash}p{#1}}

% Kanban card command
\newcommand{\kcard}[3]{%
  \par\noindent
  \colorbox{#1!15}{%
    \begin{minipage}[t]{0.95\linewidth}
      \vspace{0.3em}
      \small\textbf{\color{#1}#2}\\[0.2em]
      \footnotesize #3
      \vspace{0.3em}
    \end{minipage}%
  }\par\vspace{0.4em}
}

\begin{document}
\pagestyle{empty}

\begin{center}
{\LARGE\bfseries لوحة كانبان للمشروع البحثي}\\[0.3em]
{\large عنوان المشروع — \today}\\[1em]
\end{center}

% ============================================================
% KANBAN BOARD (4 columns)
% ============================================================
\noindent
\begin{tabularx}{\textwidth}{|X|X|X|X|}
\hline
\rowcolor{headerColor}
\textcolor{white}{\textbf{\large للمستقبل (\LR{To Do})}} &
\textcolor{white}{\textbf{\large قيد التنفيذ (\LR{In Progress})}} &
\textcolor{white}{\textbf{\large قيد المراجعة (\LR{Review})}} &
\textcolor{white}{\textbf{\large مكتمل (\LR{Done})}} \\
\hline
\end{tabularx}

\vspace{0.5em}

\noindent
\begin{tabularx}{\textwidth}{X X X X}
% ---- Column 1: To Do ----
\cellcolor{todoColor}
\kcard{todoColor}{T1: مراجعة الأدبيات}{المسؤول: Co1 \\ الأولوية: عالية}
\kcard{todoColor}{T2: تحديد الفجوات}{المسؤول: PI \\ الأولوية: عالية}
\kcard{todoColor}{T3: تصميم المنهجية}{المسؤول: PI \\ الأولوية: متوسطة}
\kcard{todoColor}{T4: تجهيز المعدات}{المسؤول: Tech \\ الأولوية: منخفضة}
&
% ---- Column 2: In Progress ----
\cellcolor{progressColor!10}
\kcard{progressColor}{P1: جمع البيانات}{المسؤول: Co1, Co2 \\ التقدم: 60\%}
\kcard{progressColor}{P2: معالجة البيانات}{المسؤول: Co2 \\ التقدم: 30\%}
&
% ---- Column 3: Review ----
\cellcolor{reviewColor!10}
\kcard{reviewColor}{R1: تقرير الأدبيات}{المراجع: PI \\ بانتظار المراجعة}
\kcard{reviewColor}{R2: تصميم التجربة}{المراجع: PI \\ بانتظار المراجعة}
&
% ---- Column 4: Done ----
\cellcolor{doneColor!10}
\kcard{doneColor}{D1: ميثاق المشروع}{أُنجز في: ---}
\kcard{doneColor}{D2: الموافقة على التمويل}{أُنجز في: ---}
\end{tabularx}

\vspace{1.5em}

% ============================================================
% WIP LIMITS & RULES
% ============================================================
\section*{حدود العمل الجاري (WIP Limits) والقواعد}

\begin{tabularx}{\textwidth}{L{4cm} X}
\toprule
\textbf{قيد التنفيذ} & حد أقصى: \textbf{3 مهام} في آن واحد \\
\textbf{قيد المراجعة} & حد أقصى: \textbf{2 مهمة} في آن واحد \\
\bottomrule
\end{tabularx}

\vspace{1em}

\textbf{قواعد لوحة كانبان:}
\begin{itemize}
    \item تُسحب البطاقة من اليمين إلى اليسار عند اكتمالها.
    \item لا يمكن تجاوز حدود العمل الجاري (\LR{WIP Limits}).
    \item تُحدّث حالة البطاقة يومياً.
    \item تُذكر العقبات في قسم القضايا (\LR{Issue Log}).
    \item تُرتّب البطاقات حسب الأولوية داخل كل عمود.
\end{itemize}

\end{document}
